%% --------------------------------------------------------
\section{Теорія збурень Møller–Plesset (MP2)}
%% --------------------------------------------------------

%% --------------------------------------------------------
\subsection{Теоретичні основи MP2}
%% --------------------------------------------------------

Теорія збурень Møller–Plesset (MP) --- один із найпоширеніших
пост-Хартрі–Фоківських методів для врахування електронної кореляції.
Її ідея полягає у розкладі точного гамільтоніана системи у вигляді суми
незбуреної частини (розв’язаної в методі Хартрі–Фока) і малих поправок:

\begin{equation}
    \hat{H} = \hat{H}_0 + \lambda \hat{V},
\end{equation}
де
\(\hat{H}_0\) --- оператор Фока,
а \(\hat{V} = \hat{H} - \hat{H}_0\) --- «збурення», яке враховує
електронну кореляцію, відсутню в одноелектронному описі HF.

Енергію можна подати у вигляді ряду за параметром збурення~$\lambda$:
\begin{equation}
    E = E^{(0)} + E^{(1)} + E^{(2)} + E^{(3)} + \ldots
\end{equation}

де:
\begin{itemize}
    \item $E^{(0)}$ --- енергія, яку дає саме наближення Хартрі–Фока;
    \item $E^{(1)}$ --- перша поправка збурень, яка, як виявляється,
          \emph{тотожно дорівнює нулю} у випадку HF (через варіаційний принцип);
    \item $E^{(2)}$ --- перша ненульова корекція, яку називають
          \textbf{енергією кореляції MP2};
    \item $E^{(3)}, E^{(4)}$ --- наступні наближення (\texttt{MP3}, \texttt{MP4}),
          що уточнюють опис кореляції.
\end{itemize}

Таким чином, \textbf{MP1 не існує як окремий метод}, оскільки
поправка першого порядку вже врахована у розв’язку рівнянь Фока:
\[
    E^{(0)} + E^{(1)} = E_{\text{HF}}.
\]

%% --------------------------------------------------------
\subsection{Формула енергії MP2}
%% --------------------------------------------------------

У другому порядку теорії збурень енергія записується як:

\begin{equation}\label{eq:MP2energy}
E^{(2)} =
\frac{1}{4}
\sum_{ijab}
\frac{|\langle ij || ab \rangle|^2}
{\varepsilon_i + \varepsilon_j - \varepsilon_a - \varepsilon_b},
\end{equation}

де:
\begin{itemize}
    \item $i,j$ --- індекси зайнятих (occupied) орбіталей,
    \item $a,b$ --- індекси віртуальних (unoccupied) орбіталей,
    \item $\langle ij || ab \rangle$ --- антисиметризовані двоелектронні інтеграли,
    \item $\varepsilon_k$ --- орбітальні енергії з HF-розрахунку.
\end{itemize}

Цей вираз описує зниження енергії системи за рахунок корельованого руху
електронів, коли збурення дозволяє їм «обмінюватися» з віртуальними станами.

%% --------------------------------------------------------
\subsection{Фізичний зміст MP2-корекції}
%% --------------------------------------------------------

\texttt{MP2} описує \emph{динамічну електронну кореляцію} —
швидкі короткодіючі флуктуації положень електронів.
Інтуїтивно, електрони «відчувають» один одного і миттєво уникають,
що знижує енергію системи відносно незалежного HF опису.

\texttt{MP2} добре працює для:
\begin{itemize}
    \item молекул поблизу рівноважної геометрії;
    \item сполук із замкненими оболонками;
    \item систем, де кореляція не надто сильна.
\end{itemize}

А от при розриві зв’язків або в багатоконфігураційних системах
(де один детермінант вже неадекватний) MP2 дає хибні результати —
у таких випадках потрібні методи на кшталт \texttt{CASSCF} або \texttt{CCSD(T)}.

%% --------------------------------------------------------
\subsection{Практичне застосування}
%% --------------------------------------------------------

У більшості сучасних програм (Gaussian, ORCA, PySCF) MP2 є
базовим рівнем пост-HF розрахунків. Він забезпечує добрий баланс
між точністю та обчислювальною вартістю, додаючи лише \(O(N^5)\)
складність порівняно з HF (\(O(N^4)\)).

\begin{center}
\begin{tabular}{lcc}
\toprule
Метод & Точність & Обчислювальна складність \\
\midrule
HF & $\sim 1-5$ kcal/mol & $O(N^4)$ \\
MP2 & $\sim 0.5-1$ kcal/mol & $O(N^5)$ \\
CCSD & $\sim 0.1$ kcal/mol & $O(N^6)$ \\
CCSD(T) & $\sim 0.01$ kcal/mol & $O(N^7)$ \\
\bottomrule
\end{tabular}
\end{center}
\texttt{}
Отже, \texttt{MP2} є першим кроком за межі одноелектронної теорії,
який дозволяє описати тонку структуру енергій, дисперсійні взаємодії
та слабкі зв’язки, недосяжні в рамках Хартрі–Фока.


%% --------------------------------------------------------
\subsection{Практичний гайд: розрахунок MP2 у PySCF}
%% --------------------------------------------------------

Метод \texttt{MP2} у PySCF реалізований як надбудова над
звичайним розрахунком Хартрі–Фока (\texttt{RHF} або \texttt{UHF}).
Тобто спочатку потрібно виконати HF-розрахунок, а потім
передати отриманий об’єкт у модуль \texttt{mp}.


\inputcode[mode=inline]{mp2_example.py}

\textbf{Пояснення кроків:}
\begin{enumerate}
    \item \inlinecode{mol} --- створення об'єкта молекули (як завжди у PySCF);
    \item \inlinecode{scf.RHF(mol)} --- запуск Хартрі–Фока;
    \item \inlinecode{mp.MP2(mf)} --- ініціалізація MP2 з уже готовим SCF-об’єктом;
    \item \inlinecode{kernel()} --- обчислення повної та кореляційної енергії.
\end{enumerate}

\textbf{Типовий вивід:}
\begin{minted}{text}
Hartree–Fock energy = -76.0267656731 Hartree
MP2 total energy      = -76.2307856559 Hartree
MP2 correlation energy = -0.2040199828 Hartree
\end{minted}

\textbf{Важливо:} Для відкритих оболонок використовуйте \inlinecode{mp.UMP2} (на базі \texttt{UHF}).

%% --------------------------------------------------------
\subsection{Практичні аспекти та обчислювальна складність MP2}
%% --------------------------------------------------------

\paragraph{Заморожене ядро (frozen core).}

В теорії MP2 енергія кореляції обчислюється за формулою~\eqref{eq:MP2energy}.

Глибокі внутрішні орбіталі (\emph{core orbitals}), наприклад \(1s\)-орбіталі кисню у молекулі води,
мають великі за модулем енергії~\(\varepsilon_i\).
Через це знаменник у наведеному виразі стає дуже великим, а внесок таких орбіталей у суму --- незначним:
\[
    \frac{|\langle ij || ab \rangle|^2}
    {\varepsilon_i + \varepsilon_j - \varepsilon_a - \varepsilon_b} \approx 0.
\]
До того ж, ці орбіталі локалізовані поблизу ядра і слабо взаємодіють з валентними електронами,
тобто практично не впливають на хімічні властивості.

Щоб уникнути зайвих обчислень, у програмі \texttt{PySCF} часто задають параметр
\inlinecode{frozen}, який «заморожує» кілька найнижчих (core) орбіталей:
\begin{minted}{python}
mp2_calc = mp.MP2(mf).set(frozen=2)
\end{minted}
Це означає, що дві найнижчі орбіталі не включаються у розрахунок енергії кореляції.
Таким чином зменшується кількість зайнятих орбіталей \(N_\text{occ}\),
а відповідно --- кількість комбінацій у сумі за \(i,j,a,b\),
що суттєво скорочує час розрахунку без втрати точності для валентної області.

Такий підхід має сенс, якщо система не містить сильно корельованих внутрішніх електронів,
тобто коли валентна і внутрішня області добре розділені за енергією.

\paragraph{Обчислювальна складність \(O(N^5)\).}

Метод \texttt{MP2} є першим після Хартрі–Фока, який враховує електронну кореляцію,
але навіть він є досить ресурсоємним.
Це зумовлено структурою формули для \(E_{\text{MP2}}\), де присутня четверна сума:
\[
    \sum_{ijab} \quad \Rightarrow \quad
    N_\text{occ}^2 N_\text{virt}^2 \sim O(N^4),
\]
де \(N_\text{occ}\) --- число зайнятих, а \(N_\text{virt}\) --- число віртуальних орбіталей.

Однак найважчий етап --- не сама сума, а перетворення двоелектронних інтегралів
з базису атомних орбіталей (AO) у базис молекулярних орбіталей (MO):
\begin{equation}
    (ij|ab) =
    \sum_{\mu\nu\lambda\sigma}
    C_{\mu i} C_{\nu j} C_{\lambda a} C_{\sigma b}
    (\mu\nu|\lambda\sigma),
\end{equation}
де \(C_{\mu p}\) --- коефіцієнти розкладу орбіталей у базисі AO.
Кожна така трансформація масштабується як \(O(N^5)\),
оскільки для кожного набору індексів потрібно виконати сумування за чотирма базисними функціями.

Отже, загальна складність MP2 зростає як
\[
    T_{\text{MP2}} \sim O(N^5),
\]
що обмежує застосування цього методу до систем з приблизно 40--60 атомами
(залежно від базису та апаратних ресурсів).

\paragraph{Оптимізація MP2 у практиці.}

Для пришвидшення розрахунків у \texttt{PySCF} застосовують такі прийоми:
\begin{itemize}
    \item Використання опції \inlinecode{frozen} для замороження ядерних орбіталей;
    \item Застосування апроксимації RI-MP2 (Resolution of Identity), яка замінює
          чотириіндексні інтеграли на трьохіндексні, зменшуючи масштабність до \(O(N^3)\);
    \item Локалізація орбіталей (LMP2), що дозволяє ефективно розв'язувати великі системи
          за рахунок відсікання слабких міжорбітальних взаємодій.
\end{itemize}

\inputcode[mode=inline]{mp2_optimization_demo.py}


\subsection{Приклади MP2-розрахунків у PySCF}
\begin{enumerate}
\item MP2 розрахунок атома Неону: \inputcode{Ne_mp2.py}
\item MP2 для відкритих оболонок (UMP2): \inputcode{mp2_open_shell.py}
\item Систематичне порівняння HF vs MP2: \inputcode{HF_vs_MP2.py}
\item Залежність \texttt{MP2} від базисного набору: \inputcode{mp2_basis_convergence.py}

\texttt{MP2} сильно залежить від базису, особливо потрібні функції високих $l$.
\end{enumerate}

%% --------------------------------------------------------
\subsection{Переваги та недоліки MP2}
%% --------------------------------------------------------

\paragraph{Переваги:}
\begin{itemize}
    \item Відносно швидкий ($\mathcal{O}(N^5)$) порівняно з іншими post-HF
    \item Добре відновлює динамічну кореляцію
    \item Size-consistent (правильне масштабування для великих систем)
    \item Доступний для великих атомів/молекул
    \item Добрий базис для вищих методів (MP3, MP4)
\end{itemize}

\paragraph{Недоліки:}
\begin{itemize}
    \item Не виправляє забруднення спіном UHF
    \item Може розходитись для систем з малим gap HOMO-LUMO
    \item Погано працює для сильно корельованих систем
    \item Потребує великих базисів для точних результатів
    \item Не варіаційний (може давати енергію нижчу за точну)
\end{itemize}

%% --------------------------------------------------------
\subsection{Рекомендації по використанню MP2}
%% --------------------------------------------------------

\begin{table}[h]
    \centering
    \caption{Коли використовувати MP2}
    \label{tab:mp2_recommendations}
    \small
    \begin{tabular}{ll}
        \hline
        \textbf{Добре для:}     & \textbf{Погано для:}          \\
        \hline
        Замкнені оболонки       & Системи з малим HOMO-LUMO gap \\
        Якісні оцінки кореляції & Розрив зв'язків               \\
        Великі системи          & Збуджені стани                \\
        Слабкі взаємодії        & Перехідні стани               \\
        Відносні енергії        & Діелектронні системи          \\
        \hline
    \end{tabular}
\end{table}
